%%%%%%%%%%%%%%%%%%%%% ch08.tex (EN) --- monograph synthesis (not a paper clone) %%%
\chapter{Digital Twin with LLMs in Inspection Analysis}
\label{ch:llm-dt}

\abstract*{This chapter develops the decision layer of the UAV inspection pipeline: after planning, reconstruction, detection, and GIS anchoring, how can large language and vision--language models support diagnosis and maintenance decisions on digital twins? We proceed from the broad to the specialised: first surveying how LLMs are already used across building inspection and facility management; then explaining the professional requirements that separate chatbots from audit-ready systems; and finally organising a twin-grounded framework of observation records, asset-centric knowledge, topology-aware retrieval, agents, and report delivery. Community work is synthesised alongside group contributions; asset-centric RAG appears only as one worked example among published siblings.}

\abstract{Chapters~2--7 explain how to fly, reconstruct, detect, and map. Chapter~8 asks what happens next when inspectors and facility managers want natural-language answers that remain engineer-auditable. The narrative follows a monograph funnel---\emph{broad applications in inspection}, \emph{professional constraints}, \emph{model foundations and grounding}, \emph{twin knowledge packaging}, \emph{multimodal retrieval and agents}, and \emph{reports plus decision metrics}. Both LLMs and VLMs are in scope. The chapter is a synthesis for the book pipeline, not a restatement of any single manuscript.}

%%%%%%%%%%%%%%%%%%%%%%%%%%%%%%%%%%%%%%%%%%%%%%%%%%%%%%%%%%%%%%%%%%%%%%
\section{Digital Twins for Building Inspection}
\label{sec:ch8-dt}

\subsection{Where this chapter sits in the book}
The monograph is organised as one pipeline. Task and motion planning (Chapters~2--3) produce safe coverage. Reconstruction (Chapter~4) turns imagery into geometry. Datasets and perception models (Chapters~5--6) name defects. GIS and GeoBIM (Chapter~7) give assets addresses in the city. None of those chapters, alone, answers questions such as ``which panels participate in a multi-storey moisture path?'' or ``draft a repair brief that cites evidence an engineer can audit.'' That is the job of the \emph{decision layer}: connecting twin entities to language, retrieval, and---increasingly---tool-using agents.

\begin{figure}[t]
\centering
\includegraphics[width=0.9\textwidth]{framework_overview_overall.pdf}
\caption{Book-level view of inspection intelligence. This chapter focuses on the language, retrieval, and decision layers that consume twin and perception products from earlier chapters.}
\label{fig:ch8-framework}
\end{figure}

\subsection{Opening scenario: questions a pure geometry twin cannot answer}
Consider a logistics warehouse after a quarterly UAV survey. An FM manager asks three questions that sound ordinary but are methodologically demanding:
\begin{enumerate}[leftmargin=*]
\item \emph{Spatial.} ``Show me every Asset-ID on elevation E with delamination severity above medium.''
\item \emph{Multi-hop / propagation.} ``If moisture appears two storeys above panel A-07, which neighbours should we inspect next?''
\item \emph{Aggregation / planning.} ``Draft a gondola schedule that prioritises assets with both RGB cracking and IR anomalies, citing evidence I can open in the twin.''
\end{enumerate}
A mesh viewer can colour faces. A detector can draw boxes. Neither produces an auditable, natural-language answer that binds symptoms to Asset-IDs, neighbours, and repair priority. The rest of this chapter builds the missing layer.

\subsection{What a twin must provide for language use}
A geometric twin answers location. A decision-ready twin must also expose (i)~stable asset identifiers, (ii)~provenance for each observation (sensor, time, detector), and (iii)~relations that support propagation reasoning. Without those, language models can only chat over loose files---fluent, but not auditable for facility management (FM).

Digital-twin surveys in construction emphasise that twins are not merely 3D viewers: they are evolving information systems linking as-built geometry, semantics, and operations~\cite{boje2020,opoku2021,lu2020}. UAV--BIM inspection reviews sketched early automation agendas~\cite{zhang2022}. Within this research group, Zhang et al.\ demonstrated AI-empowered twin modelling that couples UAV sensing with GeoBIM for high-precision defect management~\cite{zhang2025}; Det-Recon-Reg packages detection, reconstruction, and registration for large-scale inspection~\cite{yang2025detreconreg}. Related twin uses extend beyond facades, including worker monitoring on large sites~\cite{han2025}. Multi-agent virtual modelling inside BIM further shows that twins are becoming active information systems, not static viewers~\cite{lee2025multiagentvbm}. These works supply the geometric and identifier premise for the language layer developed below.

\begin{figure}[t]
\centering
\includegraphics[width=0.88\textwidth]{dt_modeling_pipeline.pdf}
\caption{Typical UAV-to-twin modelling path that earlier chapters develop in depth. Here it is treated as the \emph{input interface} to language-level analysis.}
\label{fig:ch8-dt-pipeline}
\end{figure}

\subsection{Gaps that language alone cannot fill}
Three gaps recur in field programmes:
\begin{itemize}[leftmargin=*]
\item \textbf{Visualisation $\neq$ reasoning.} Seeing a crack does not explain cross-storey risk.
\item \textbf{Missing provenance.} Without sensor/time/detector stamps, generated advice cannot be audited.
\item \textbf{Missing natural-language interface with discipline.} Chat over PDFs invents Asset-IDs; chat over twins without citation rules does the same more confidently.
\end{itemize}
Chapter~8 therefore promises structured knowledge units, grounded LLM/VLM generation, constrained reports, and decision metrics---not a chatbot glued onto a point cloud.

\paragraph{Bridge.}
We therefore begin with the \emph{broad} landscape of how large models are already used in inspection and facility practice, then narrow to the \emph{specialised} requirements of twin-grounded facade diagnosis.

%%%%%%%%%%%%%%%%%%%%%%%%%%%%%%%%%%%%%%%%%%%%%%%%%%%%%%%%%%%%%%%%%%%%%%
\section{Foundations of LLMs for Civil Engineering Applications}
\label{sec:ch8-llm-foundations}

\subsection{Broad view: where large models already help inspection}
Large language models (LLMs) and vision--language models (VLMs) have entered building inspection not as a single product, but as a family of interfaces between humans, documents, models, and machines~\cite{saka2024gptconstruction,ma2026tenquestions}. At the broadest level, practitioners use them for five recurring jobs:

\begin{enumerate}[leftmargin=*]
\item \textbf{Document and code question answering.} Staff ask about editions of codes, safety manuals, and O\&M handbooks; retrieval-augmented generation (RAG) typically outperforms relying on parametric memory alone~\cite{joffe2025buildingcodes,lee2024ragsafety,krutli2025fmrags}.
\item \textbf{BIM and model navigation.} Natural-language assistants locate objects, properties, and relations inside design or as-built BIM~\cite{zheng2023bimassistant,guo2025bimllmquery}.
\item \textbf{Defect dialogue and triage.} Hybrid RAG supports conversations over defect records and management logs~\cite{jeon2025hybridrag}, drafting candidate priorities for human review.
\item \textbf{Report drafting and summarisation.} Models turn inspection notes, photo captions, and checklist outcomes into first-draft narratives that engineers then edit.
\item \textbf{Task brokering toward twins and robots.} Emerging agents translate user intent into tool calls that read a twin, request a robot photo, or stage an update after approval~\cite{luo2025llmdtrobot,gao2026ifcagent,rawat2025llmmas,lee2025multiagentvbm}.
\end{enumerate}

\begin{table}[t]
\centering
\caption{Broad application modes of LLMs/VLMs in building inspection and FM (landscape before specialisation).}
\label{tab:ch8-llm-modes}
\begin{tabular}{p{2.6cm}p{3.6cm}p{5.4cm}}
\toprule
Mode & Typical input & Typical output \\
\midrule
Code / manual QA & regulation text, handbooks & cited clauses, checklists \\
BIM navigation & IFC/BIM objects, queries & object-linked answers \\
Defect dialogue & logs, images, prior tickets & triage notes, hypotheses \\
Report drafting & notes + photo captions & editable report draft \\
Twin / robot agency & twin state + user intent & gated tool calls / tasks \\
\bottomrule
\end{tabular}
\end{table}

These modes already matter for UAV programmes: Chapters~5--6 produce observations; Chapters~4 and~7 produce geometry and Asset-IDs; facility teams still need language interfaces to manuals, BIM, and work orders. Domain primers emphasise both opportunity and reliability limits~\cite{ma2026tenquestions,saka2024gptconstruction}. The remainder of this section first clarifies \emph{why} generic chat is insufficient for professional inspection, then reviews the model foundations and grounding mechanisms required before we specialise to twin-grounded facade reasoning in Sections~\ref{sec:ch8-dt-to-language}--\ref{sec:ch8-dss}.

\subsection{From broad use to professional inspection requirements}
Inspection and FM decisions are consequential: wrong Asset-IDs waste gondola time; unauditable advice fails compliance review; silent averaging of RGB/IR conflict can hide moisture risk. Relative to open-domain chat, professional inspection imposes four constraints:

\begin{itemize}[leftmargin=*]
\item \textbf{Identity discipline.} Answers must bind claims to stable Asset-IDs shared with GIS/FM systems (Chapter~7), not to invented labels.
\item \textbf{Evidence provenance.} Each material claim should open an observation (sensor, time, detector, image URI) produced by Chapters~5--6 and registered by Chapters~4--7.
\item \textbf{Multi-modal honesty.} Visual and thermal cues may disagree; language systems must disclose conflict rather than flatten it~\cite{zhao2024facade,zhao2025icra}.
\item \textbf{Action gating.} Twin writes, UAV revisits, and robot dispatch are irreversible relative to chat; they require human approval and safety checks~\cite{luo2025llmdtrobot}.
\end{itemize}

Hallucination surveys explain the underlying risk: fluent generators invent entities and citations when evidence is missing or weakly retrieved~\cite{ji2023}. Post-hoc verbal justifications are not evidence grounding~\cite{zhao2024}. Hence the professional question is not ``which chatbot brand?'' but ``how is the model grounded on editioned corpora and twin entities?''

\subsection{Minimal technical foundations (what the models are)}
Contemporary LLMs rest on the Transformer architecture~\cite{vaswani2017attention}: self-attention models long-range dependencies in specifications, logs, and multi-sentence inspection questions. Large-scale pretraining yields few-shot learners~\cite{brown2020}; instruction tuning aligns models to user intent~\cite{ouyang2022instructgpt}; chain-of-thought prompting can elicit intermediate steps for multi-hop questions~\cite{wei2022cot}. None of these alone guarantees factual correctness about a particular facade.

\begin{figure}[t]
\centering
\begin{tikzpicture}[
  box/.style={draw, rounded corners=2pt, align=center, minimum height=1.05cm, minimum width=2.35cm, font=\small},
  arr/.style={-{Latex}, thick}
]
\node[box, fill=black!6] (p) {Pretrain\\[-1pt]{\scriptsize parametric memory}};
\node[box, fill=black!6, right=0.45cm of p] (i) {Instruct /\\[-1pt]{\scriptsize align to intent}};
\node[box, fill=black!6, right=0.45cm of i] (c) {CoT /\\[-1pt]{\scriptsize multi-step talk}};
\node[box, fill=black!12, right=0.45cm of c] (h) {Still fails:\\[-1pt]{\scriptsize hallucination, wrong IDs}};
\draw[arr] (p) -- (i);
\draw[arr] (i) -- (c);
\draw[arr] (c) -- (h);
\node[below=0.35cm of p, font=\scriptsize, text width=2.4cm, align=center] {knows language};
\node[below=0.35cm of i, font=\scriptsize, text width=2.4cm, align=center] {follows prompts};
\node[below=0.35cm of c, font=\scriptsize, text width=2.4cm, align=center] {explains steps};
\node[below=0.35cm of h, font=\scriptsize, text width=2.5cm, align=center] {needs grounding\\on twin evidence};
\end{tikzpicture}
\caption{Capability ladder for LLMs in inspection dialogue. Fluency arrives before auditability; grounding is a separate engineering layer.}
\label{fig:ch8-llm-ladder}
\end{figure}

Treat an LLM as a conditional text generator whose knowledge is (i)~frozen at training cut-off in parametric memory and (ii)~easily overwritten by the prompt. That dual nature is why grounding strategies matter more than brand names when answers must cite real Asset-IDs.

Inspection data are not text-only. Contrastive vision--language pretraining (e.g., CLIP) aligns images and captions~\cite{radford2021clip}; instruction-tuned multimodal models accept images with text prompts~\cite{liu2023llava}. In this book, VLMs belong to the decision layer: they describe visual evidence, compare RGB/IR cues in language, and support report drafting. They complement---rather than replace---Chapter~6 detectors when audit trails require typed labels, confidences, and host IDs.

\begin{figure}[t]
\centering
\begin{tikzpicture}[
  role/.style={draw, rounded corners=2pt, align=center, minimum height=1.35cm, minimum width=3.4cm, font=\small},
  arr/.style={-{Latex}, thick}
]
\node[role, fill=black!5] (det) {\textbf{Detector / VLM head}\\writes typed observations};
\node[role, fill=black!5, right=0.7cm of det] (twin) {\textbf{Twin / GIS}\\Asset-ID + topology};
\node[role, fill=black!10, right=0.7cm of twin] (llm) {\textbf{LLM / VLM}\\retrieve, narrate, propose};
\draw[arr] (det) -- (twin);
\draw[arr] (twin) -- (llm);
\node[below=0.25cm of det, font=\scriptsize] {Ch.~5--6};
\node[below=0.25cm of twin, font=\scriptsize] {Ch.~4, 7};
\node[below=0.25cm of llm, font=\scriptsize] {Ch.~8};
\end{tikzpicture}
\caption{Role split among perception, twin/GIS, and language stacks for professional inspection.}
\label{fig:ch8-vlm-roles}
\end{figure}

\subsection{Memory and grounding: the professional control stack}
Before specialising to facade twins, separate four memory regimes:

\begin{enumerate}[leftmargin=*]
\item \textbf{Parametric memory} --- facts in weights; hard to update after every flight~\cite{brown2020}.
\item \textbf{Working / contextual memory} --- prompt window and short dialogue state.
\item \textbf{Non-parametric retrieval memory} --- RAG indexes over editioned manuals or observation cards~\cite{lewis2020,shuster2021,krutli2025fmrags}.
\item \textbf{Structured long-term twin memory} --- Asset-IDs, topology, and time-stamped observations in BIM/GIS/twin stores~\cite{pan2024b,zhang2025,park2023generativeagents,rawat2025llmmas}.
\end{enumerate}

\begin{figure}[t]
\centering
\begin{tikzpicture}[
  m/.style={draw, align=center, minimum height=1.15cm, minimum width=2.55cm, font=\small, rounded corners=2pt},
  arr/.style={-{Latex}, thick}
]
\node[m, fill=black!3] (a) {Parametric\\[-1pt]{\scriptsize weights}};
\node[m, fill=black!5, right=0.35cm of a] (b) {Contextual\\[-1pt]{\scriptsize prompt/dialog}};
\node[m, fill=black!8, right=0.35cm of b] (c) {Retrieval\\[-1pt]{\scriptsize RAG index}};
\node[m, fill=black!14, right=0.35cm of c] (d) {Twin store\\[-1pt]{\scriptsize Asset-ID + topo}};
\draw[arr] (a) -- (b);
\draw[arr] (b) -- (c);
\draw[arr] (c) -- (d);
\node[below=0.4cm of a, font=\scriptsize, text width=2.5cm, align=center] {slow to update};
\node[below=0.4cm of b, font=\scriptsize, text width=2.5cm, align=center] {session-limited};
\node[below=0.4cm of c, font=\scriptsize, text width=2.5cm, align=center] {editioned docs /\\observation cards};
\node[below=0.4cm of d, font=\scriptsize, text width=2.6cm, align=center] {\textbf{authoritative}\\for facade FM};
\end{tikzpicture}
\caption{Four memory regimes. Broad document RAG uses the third; twin-grounded inspection must also master the fourth.}
\label{fig:ch8-memory-ladder}
\end{figure}

Four complementary grounding strategies appear across AEC deployments:

\begin{enumerate}[leftmargin=*]
\item \textbf{Prompting and schemas.} Constrained JSON/report templates reduce free-form drift~\cite{wei2022cot}.
\item \textbf{RAG.} Attach generation to external memory so answers can cite retrieved passages~\cite{lewis2020,shuster2021,gao2023hyde,asai2024selfrag,jeong2024adaptive,peng2024}.
\item \textbf{Tools and agents.} Reason--act loops call APIs, databases, or robots~\cite{yao2023react,schick2023toolformer,wang2024surveyagents,li2024surveymultiagent}.
\item \textbf{Structured knowledge.} Typed entities and graphs improve constraint checking~\cite{pan2024b}---for buildings: Asset-IDs, IFC objects, observation records.
\end{enumerate}

\begin{figure}[t]
\centering
\begin{tikzpicture}[
  g/.style={draw, rounded corners=2pt, align=center, minimum height=1.0cm, minimum width=2.4cm, font=\scriptsize},
  arr/.style={-{Latex}, thick}
]
\node[g, fill=black!3] (s1) {Prompt /\\schema};
\node[g, fill=black!6, right=0.4cm of s1] (s2) {RAG\\retrieval};
\node[g, fill=black!9, right=0.4cm of s2] (s3) {Tools /\\agents};
\node[g, fill=black!14, right=0.4cm of s3] (s4) {Structured\\Asset-IDs};
\draw[arr] (s1) -- (s2);
\draw[arr] (s2) -- (s3);
\draw[arr] (s3) -- (s4);
\node[below=0.55cm of $(s1)!0.5!(s4)$, font=\small] {Stronger auditability $\longrightarrow$};
\end{tikzpicture}
\caption{Grounding stack from broad AEC practice to specialised twin discipline.}
\label{fig:ch8-grounding-stack}
\end{figure}

\subsection{Specialising to UAV facade inspection on digital twins}
Broad modes above are necessary but not sufficient for the questions in Section~\ref{sec:ch8-dt}. Code RAG answers ``what does the edition say?''; BIM assistants answer ``where is this object?''; facade O\&M further asks ``which maintainable assets participate in a propagation path, and can every claim open twin evidence?'' That specialisation has three design consequences:

\begin{enumerate}[leftmargin=*]
\item \textbf{Retrieve over assets and neighbours}, not only document chunks---facility questions are asset-centric, and moisture/crack paths need topology from BIM/GIS adjacency~\cite{peng2024,jeon2025hybridrag}.
\item \textbf{Cite-or-abstain on twin IDs.} Every Asset-ID / observation ID in the answer must appear in the retrieved set; otherwise abstain or escalate~\cite{ji2023,shuster2021}.
\item \textbf{Keep expert corpora and twin observations in separate stores.} Editioned manuals remain valuable~\cite{krutli2025fmrags,joffe2025buildingcodes}, but mixing them with observation cards invites the model to ``repair'' a missing Asset-ID with a generic handbook clause.
\end{enumerate}

Cognitive-twin and agent programmes show the same trajectory from answering toward acting on interoperable knowledge~\cite{shi2025cognitivedt,luo2025llmdtrobot,gao2026ifcagent,zhang2025llmagentschema}. Perception context still matters: language systems inherit the granularity and bias of facade detectors and quantification models~\cite{zhao2024facade,zhao2025icra,zhao2026physical}. The next sections therefore package twin products as observation records and asset passports, assemble multimodal topology-aware context, and only then discuss agents, reports, and decision metrics.

\paragraph{Bridge.}
Having moved from broad inspection uses to specialised twin requirements, we now specify how Chapters~4--7 products should be packaged so grounding strategies have reliable objects to retrieve.

%%%%%%%%%%%%%%%%%%%%%%%%%%%%%%%%%%%%%%%%%%%%%%%%%%%%%%%%%%%%%%%%%%%%%%
\section{From Digital Twins to Language Reasoning}
\label{sec:ch8-dt-to-language}

\subsection{A specialised teaching framework (after the broad landscape)}
Section~\ref{sec:ch8-llm-foundations} surveyed broad inspection uses of LLMs. Here we specialise: rather than importing any one paper's module list, we keep three layers separate for twin-grounded facade reasoning:

\begin{enumerate}[leftmargin=*]
\item \textbf{Observation records} --- one defect sighting with sensor provenance;
\item \textbf{Asset passports} --- all observations and geometry/topology for a maintainable host;
\item \textbf{Query assembly} --- how a natural-language question selects passports and neighbours before generation.
\end{enumerate}

Different research groups implement these layers with different schemas. The pedagogical point is the \emph{separation}: perception writes observations; FM thinks in assets; language models must not be asked to invent either. This stance is consistent with structured-knowledge + LLM roadmaps~\cite{pan2024b} and with BIM-IR systems that refuse to treat the model as a database~\cite{guo2025bimllmquery,zheng2023bimassistant}.

\subsection{Observation records in practice}
Across UAV inspection programmes, useful observation records typically store: host asset ID, time, modality (RGB/IR/\ldots), detector signature, pose or surface coordinates, and paths to imagery. GeoBIM-DT pipelines show how registration attaches detections to BIM hosts before any LLM is involved~\cite{zhang2025,yang2025detreconreg}. Chapter~6 supplies the detector; Chapters~4--7 supply the host ID.

For teaching, one may write a generic observation as a four-part record
\begin{equation}
o = \langle I,\, E,\, S,\, A \rangle,
\label{eq:ch8-obs}
\end{equation}
where $I$ collects imagery pointers, $E$ evidence fields (detector, confidence, time), $S$ semantic labels (defect type, severity hypothesis), and $A$ the host Asset-ID after registration. Natural-language location phrases (``top centre of the southeast panel'') help LLMs, while metric coordinates remain for field navigation. Provenance is not optional metadata: without $E$, cite-or-abstain rules have nothing to check.

\begin{figure}[t]
\centering
\includegraphics[width=0.85\textwidth]{idp_construction_pipeline.jpg}
\caption{Example pipeline that turns registered detections into structured observation records (illustrative figure from a UAV inspection stack). Read as ``how Chapters~4--6 feed Chapter~8'', not as a mandatory schema name.}
\label{fig:ch8-idp}
\end{figure}

\subsection{Asset passports and topology}
Facility questions are usually asset-centric. Passports therefore aggregate observations per maintainable host and attach (i)~symptom summaries, (ii)~host geometry for severity normalisation, and (iii)~neighbour links (above/below/lateral). Neighbour links enable propagation-style reading---e.g., mid-storey delamination considered together with upstairs moisture---without pretending a document-similarity graph is a structural graph. Graph-oriented RAG~\cite{peng2024} is related in spirit; the inspection-specific requirement is that graph edges come from BIM/GIS adjacency and registration, not only from text co-occurrence.

Teaching notation for a passport of host $A$ is
\begin{equation}
P(A) = \bigl\{\, o \mid A(o)=A \,\bigr\}
\;\cup\;
G(A)
\;\cup\;
N_k(A),
\label{eq:ch8-passport}
\end{equation}
where $G(A)$ is host geometry useful for normalisation, and $N_k(A)$ is the $k$-hop neighbour set under BIM/GIS adjacency. The hop radius $k$ is an \emph{operational} knob: moisture and vertical crack paths often need $k\ge 1$; pure local lookups may use $k=0$.

\subsection{A five-layer index for language queries}
Vector search alone answers ``which text chunks look similar?'' Facade FM questions need more. A teaching index stack is:

\begin{table}[t]
\centering
\caption{Five index layers that answer different FM questions.}
\label{tab:ch8-index}
\begin{tabular}{clp{7.2cm}}
\toprule
Layer & Answers & Typical keys \\
\midrule
L1 Identity & Which Asset-ID? & stable host ID, IFC GUID \\
L2 Spatial & Where on the facade? & elevation, storey, UV/pose \\
L3 Symptom & What was observed? & defect type, severity, modality \\
L4 Topology & Who are neighbours? & above/below/lateral edges \\
L5 Evidence & Can I open the proof? & image URI, detector, time \\
\bottomrule
\end{tabular}
\end{table}

Dense embeddings remain useful for free-text symptom phrases and manual clauses; they should sit \emph{beside} L1--L5, not replace them.

\begin{table}[t]
\centering
\caption{How earlier chapters feed the decision layer (conceptual checklist).}
\label{tab:ch8-interfaces}
\begin{tabular}{llp{5.8cm}}
\toprule
Book chapter & Product consumed here & Why language needs it \\
\midrule
Ch.~4 Reconstruction & geometry / BIM hosts & stable places to attach defects \\
Ch.~5--6 Perception & labelled observations & typed symptoms + confidence \\
Ch.~7 GIS / GeoBIM & city-scale Asset-IDs & shared identifiers with FM systems \\
This chapter & retrieval + agents + reports & traceable answers and actions \\
\bottomrule
\end{tabular}
\end{table}

\subsection{Worked example: asset-centric retrieval}
One concrete implementation pattern---used in ongoing group experiments and therefore \emph{not} entered here as a bibliography item until peer-published---organises observation records into asset passports and expands neighbours along BIM adjacency for multi-hop diagnosis. In monograph voice we treat it as a teaching pattern among many published siblings: hybrid RAG for defect dialogue~\cite{jeon2025hybridrag}, code/safety RAG~\cite{joffe2025buildingcodes,lee2024ragsafety}, facility-document RAG~\cite{krutli2025fmrags}, and BIM assistants~\cite{zheng2023bimassistant}. Readers who need implementation minutiae should consult those primary sources; this chapter keeps the ideas needed to connect the book pipeline to decision support.

\paragraph{Bridge.}
We next assemble evidence under multimodal and topological constraints during a single query.

%%%%%%%%%%%%%%%%%%%%%%%%%%%%%%%%%%%%%%%%%%%%%%%%%%%%%%%%%%%%%%%%%%%%%%
\section{Multi-Modal Fusion and Topology-Aware RAG}
\label{sec:ch8-multimodal}

\subsection{Pipeline fusion versus end-to-end VLMs}
Facade programmes are multi-modal, and VLMs are a real part of the decision layer~\cite{radford2021clip,liu2023llava}. ``Fusion'', however, need not collapse into one opaque end-to-end VLM for every query. A book-consistent pattern that remains auditable today is: detectors (or a VLM perception head) write modality-tagged observations; retrieval selects asset passports; an LLM or VLM then composes answers under citation rules. End-to-end VLM-native twins that jointly read imagery and twin topology remain an active research direction (Chapter~9); the pedagogical point is to use VLMs where they add visual grounding, while keeping Asset-IDs and provenance outside the model's free text.

\begin{figure}[t]
\centering
\includegraphics[width=0.9\textwidth]{rag_workflow_topology.pdf}
\caption{Conceptual retrieval flow: parse intent, fetch asset passports, expand neighbours, then generate with citations. Teaching schematic for topology-aware assembly.}
\label{fig:ch8-rag}
\end{figure}

\subsection{Three-phase query assembly}
A practical query can be taught as three phases:

\begin{enumerate}[leftmargin=*]
\item \textbf{Intent parse.} Classify the question as local lookup, asset-health summary, propagation/root-cause, or floor-/elevation-level audit. Intent should change \emph{which neighbours enter the context}, not only the prompt wording.
\item \textbf{Retrieve and expand.} Fetch $P(A)$ for seed assets, then form
\[
\text{context}
=
P(A)
\cup
N_k(A)
\cup
M_{\text{opt}},
\]
where $M_{\text{opt}}$ are optional edition-versioned manuals/codes~\cite{krutli2025fmrags,joffe2025buildingcodes}.
\item \textbf{Generate under cite-or-abstain.} Every cited asset/defect ID in the answer must appear in the retrieved set; otherwise abstain or escalate to human review~\cite{lewis2020,ji2023}.
\end{enumerate}

RAG and graph-RAG surveys emphasise that retrieval quality dominates generation quality when facts are corpus-bound~\cite{lewis2020,peng2024}. Adaptive and reflective RAG variants~\cite{asai2024selfrag,jeong2024adaptive} can rewrite queries or deepen retrieval; for inspection they remain useful only if Asset-ID constraints stay hard.

\subsection{RGB/IR disagreements}
When thermal and visual evidence disagree, good practice is disclosure rather than silent averaging: RGB-only evidence should not be escalated to subsurface claims; IR-only evidence should keep uncertainty flags. Perception chapters define what each modality can see; this chapter defines how language systems must talk about conflict. VLMs may help \emph{describe} the disagreement; they must not erase the disagreement in fluent prose.

\subsection{Expert corpora versus twin observations}
Codes, repair manuals, and safety guides are valuable~\cite{lee2024ragsafety,joffe2025buildingcodes}, but they answer different questions from twin observations. Mixing them in one undifferentiated index invites the model to ``repair'' a missing Asset-ID with a generic clause from a handbook. Keep two stores: (i)~twin observation/passport memory; (ii)~editioned expert memory. Cite which store each sentence used.

\subsection{From static RAG to bounded agency}
Frontier systems increasingly interleave retrieval with tool calls (agentic RAG). For inspection we recommend \emph{bounded agency}: multi-step retrieval is welcome; final claims remain citation-locked to twin entities~\cite{yao2023react,asai2024selfrag}. Unbounded tool loops without Asset-ID checks recreate hallucination with extra latency~\cite{ji2023,shuster2021}.

\paragraph{Bridge.}
Agents become essential when the system must \emph{act}, not only answer.

%%%%%%%%%%%%%%%%%%%%%%%%%%%%%%%%%%%%%%%%%%%%%%%%%%%%%%%%%%%%%%%%%%%%%%
\section{Agentic Workflows for Twin-Grounded Inspection}
\label{sec:ch8-agents}

\subsection{Why the book needs an agent section}
RAG answers evidence-seeking questions. FM operations also require procedures: open a ticket, propose a UAV revisit, request a robot photo, or stage a twin update after human approval. That is the scope of agents~\cite{yao2023react,wang2024surveyagents,park2023generativeagents}. Building-application primers likewise list tool use and process control as open challenges~\cite{ma2026tenquestions}.

\subsection{What others have built}
Luo et al.\ place an LLM between the user, the twin, and robot middleware, emphasising human-approved twin commits~\cite{luo2025llmdtrobot}. Gao et al.\ show schema-guided multi-agent tool composition on IFC~\cite{gao2026ifcagent}. Rawat et al.\ propose construction LLM multi-agent systems with hybrid graph/vector memory~\cite{rawat2025llmmas}. Lee et al.\ demonstrate multi-agent virtual in-situ modelling of digital twins inside BIM~\cite{lee2025multiagentvbm}. Energy-oriented agent libraries show how domain tools can be versioned~\cite{zhang2025llmagentschema}. None of these replace Chapters~2--7; they consume twin state those chapters create.

\subsection{A pedagogical agent stack for facade programmes}
We describe four \emph{roles} (logical, not necessarily four model instances):
\begin{enumerate}[leftmargin=*]
\item \textbf{Retriever} --- asset-passport search with citation packing;
\item \textbf{Topology aide} --- neighbour expansion and conflict checks;
\item \textbf{Report clerk} --- fills a constrained schema (Section~\ref{sec:ch8-reports});
\item \textbf{Action proposer} --- tickets / revisit polygons / IFC patches, always behind approval.
\end{enumerate}

\begin{figure}[t]
\centering
\begin{tikzpicture}[
  r/.style={draw, rounded corners=2pt, align=center, minimum height=1.05cm, minimum width=2.2cm, font=\scriptsize},
  gate/.style={draw, thick, align=center, minimum height=1.05cm, minimum width=2.3cm, font=\scriptsize, fill=black!12},
  arr/.style={-{Latex}, thick}
]
\node[r, fill=black!3] (r1) {Retriever};
\node[r, fill=black!5, right=0.35cm of r1] (r2) {Topology\\aide};
\node[r, fill=black!7, right=0.35cm of r2] (r3) {Report\\clerk};
\node[r, fill=black!9, right=0.35cm of r3] (r4) {Action\\proposer};
\node[gate, right=0.45cm of r4] (g) {Human\\approval};
\draw[arr] (r1) -- (r2);
\draw[arr] (r2) -- (r3);
\draw[arr] (r3) -- (r4);
\draw[arr] (r4) -- (g);
\node[below=0.45cm of g, font=\scriptsize, text width=2.6cm, align=center] {twin write / robot / UAV\\only after gate};
\end{tikzpicture}
\caption{Agent role stack with a mandatory approval gate before irreversible actions.}
\label{fig:ch8-agent-stack}
\end{figure}

\begin{table}[t]
\centering
\caption{Paradigms compared for teaching (not as a bake-off among papers).}
\label{tab:ch8-paradigms}
\begin{tabular}{p{2.2cm}p{3.1cm}p{3.2cm}p{3.2cm}}
\toprule
Aspect & Chunk RAG & Asset/topology RAG & Agentic twin loop \\
\midrule
Unit & text chunk & asset passport + neighbours & passports + tools \\
State change & no & no & yes (gated) \\
Audit focus & passage cites & Asset-ID + cites & cites + approvals \\
Fits book chapters & manuals / codes & Ch.~4--7 products & Ch.~2--3 tasking too \\
\bottomrule
\end{tabular}
\end{table}

\paragraph{Micro-scenario (teaching).}
``If seepage propagates two storeys above panel A-07, which assets need gondola access?'' Retriever loads A-07; topology aide expands vertical/lateral neighbours; report clerk drafts a priority list with evidence IDs; action proposer suggests a revisit polygon for Chapters~2--3 planning---but does not dispatch until approval.

\subsection{Tool catalogue (typed surface)}
A small typed tool surface beats free-form SQL generation~\cite{gao2026ifcagent,schick2023toolformer}:

\begin{table}[t]
\centering
\caption{Minimal tool catalogue for twin-grounded inspection agents.}
\label{tab:ch8-tools}
\begin{tabular}{llp{5.6cm}}
\toprule
Tool & Side effect & Guard \\
\midrule
\texttt{get\_passport} / \texttt{search} & none & Asset-ID whitelist \\
\texttt{expand\_neighbours} & none & max hop $k$ \\
\texttt{list\_observations} & none & time window \\
\texttt{open\_evidence} & none & URI from retrieved set \\
\texttt{draft\_report} & draft only & schema validation \\
\texttt{propose\_revisit} & proposal & geofence preview \\
\texttt{propose\_twin\_patch} & proposal & human approval \\
\bottomrule
\end{tabular}
\end{table}

\subsection{Non-negotiables}
No invented Asset-IDs; no silent twin overwrites; no robot/UAV dispatch without geofence and safety checks~\cite{luo2025llmdtrobot,rawat2025llmmas}. Metrics that only score answer strings are insufficient for agent claims (Section~\ref{sec:ch8-dss}).

\paragraph{Bridge.}
Answers and proposals must land in artefacts FM systems already understand: reports.

%%%%%%%%%%%%%%%%%%%%%%%%%%%%%%%%%%%%%%%%%%%%%%%%%%%%%%%%%%%%%%%%%%%%%%
\section{Inspection Report Generation and Management}
\label{sec:ch8-reports}

\subsection{Schema-first generation}
Free-form paragraphs are hard to audit. A practical report schema groups identity (asset/defect IDs), diagnosis, propagation notes, recommended actions, and trust flags (confidence, human verification). Identity fields must be subsets of retrieved evidence. Constrained generation is the report-level analogue of structured decoding used in BIM assistants and code QA systems~\cite{zheng2023bimassistant,joffe2025buildingcodes}.

\begin{table}[t]
\centering
\caption{Example report fields for teaching constrained generation.}
\label{tab:ch8-schema}
\begin{tabular}{llp{6.0cm}}
\toprule
Group & Keys & Rule of thumb \\
\midrule
Identity & asset\_ids, defect\_ids & $\subseteq$ retrieved set \\
Diagnosis & type, severity, hypothesis & supported by observations \\
Propagation & neighbour path & from topology aide \\
Action & repair, priority & decision support, not autopilot \\
Trust & confidence, human\_check & escalate on conflict \\
\bottomrule
\end{tabular}
\end{table}

\subsection{Platform architecture as a delivery pattern}
Reports do not appear from a single prompt. A teaching architecture separates (i)~twin/GIS services that own Asset-IDs, (ii)~retrieval services that assemble passports and neighbours, (iii)~generation services that fill schemas, and (iv)~field clients that keep site tags aligned with the twin. Figure~\ref{fig:ch8-platform-arch} shows one such pattern; Figures~\ref{fig:ch8-platform}--\ref{fig:ch8-field-ui} illustrate interface consequences, not product claims.

\begin{figure}[t]
\centering
\includegraphics[width=0.88\textwidth]{platform_architecture.pdf}
\caption{Example service split for twin-grounded language analysis: identity services, retrieval/assembly, constrained generation, and field clients. Read as a delivery pattern for Section~\ref{sec:ch8-reports}.}
\label{fig:ch8-platform-arch}
\end{figure}

\begin{figure}[t]
\centering
\includegraphics[width=0.82\textwidth]{platform_demo.pdf}
\caption{Example twin UI for NL querying and report drafting. Shown as an interface pattern, not as a product claim.}
\label{fig:ch8-platform}
\end{figure}

\begin{figure}[t]
\centering
\includegraphics[width=0.88\textwidth]{multi_platform_field_interfaces.pdf}
\caption{Field interfaces that keep Asset-IDs consistent between site and twin.}
\label{fig:ch8-field-ui}
\end{figure}

\subsection{Field identity hygiene}
Language quality collapses when site tags and twin IDs diverge. Registration audits (Fig.~\ref{fig:ch8-asset-id}) are therefore part of the decision layer, not an IT afterthought---consistent with GeoBIM-era twin practice~\cite{zhang2025}.

\begin{figure}[t]
\centering
\includegraphics[width=0.85\textwidth]{field_asset_id_stages.pdf}
\caption{Staged Asset-ID verification between field capture and retrieval.}
\label{fig:ch8-asset-id}
\end{figure}

\subsection{Box~8.1: end-to-end natural-language walkthrough}
\begin{tcolorbox}[title={Box 8.1 --- From question to audit-ready draft}, breakable, colback=black!3, colframe=black!55]
\textbf{Question (FM manager).}
``Water ingress seems to climb from panel A-07 on the east elevation. Which assets two storeys above should we prioritise for gondola access, and can you draft a repair brief I can audit?''

\textbf{Step 1 --- Intent.}
Classifier marks the query as \emph{propagation + planning}. Seed asset $=\{$A-07$\}$; requested hop $k=2$; modalities of interest $=$ RGB and IR.

\textbf{Step 2 --- Retrieve.}
\texttt{get\_passport(A-07)} returns observations $o_1$ (RGB crack, conf.\ 0.91) and $o_2$ (IR anomaly, conf.\ 0.84), both registered to A-07. Manual clauses stay in a separate expert store and are not yet loaded.

\textbf{Step 3 --- Expand.}
\texttt{expand\_neighbours(A-07, $k$=2)} yields vertical chain A-07 $\rightarrow$ A-18 $\rightarrow$ A-29 and lateral neighbours A-08, A-06. Passports with moisture-related symptoms are packed into context; unrelated assets are dropped.

\textbf{Step 4 --- Conflict check.}
On A-18, RGB shows surface staining while IR is weak. Topology aide flags \texttt{modality\_conflict=true} rather than asserting subsurface leakage.

\textbf{Step 5 --- Generate (cite-or-abstain).}
Report clerk fills Table~\ref{tab:ch8-schema}:
\begin{itemize}[leftmargin=*,nosep]
\item \texttt{asset\_ids} = \{A-07, A-18, A-29\} (all $\subseteq$ retrieved set);
\item \texttt{propagation} = vertical path with evidence IDs $o_1,o_2,\ldots$;
\item \texttt{priority} = A-29 high (stacked symptoms), A-18 medium (conflict flag);
\item \texttt{human\_check} = required for A-18.
\end{itemize}
No Asset-ID outside the retrieved set is allowed. If A-29 lacked any observation, the clerk would abstain on that ID instead of inventing severity.

\textbf{Step 6 --- Action proposal (gated).}
Action proposer emits a revisit polygon hint for Chapters~2--3 and a draft gondola sequence. It does \emph{not} dispatch a UAV or write the twin until a human approves~\cite{luo2025llmdtrobot}.

\textbf{Deliverable.}
An audit-ready brief: every claim opens an evidence URI in the twin UI; conflicts are visible; irreversible actions wait behind the gate.
\end{tcolorbox}

\paragraph{Bridge.}
We close with decision-support evaluation---criteria first, case numbers second.

%%%%%%%%%%%%%%%%%%%%%%%%%%%%%%%%%%%%%%%%%%%%%%%%%%%%%%%%%%%%%%%%%%%%%%
\section{Decision Support and Predictive Maintenance}
\label{sec:ch8-dss}

\subsection{What to evaluate in a book context}
Monograph readers need criteria that travel across implementations---not a single bake-off chart:

\begin{figure}[t]
\centering
\begin{tikzpicture}[
  c/.style={draw, rounded corners=2pt, align=center, minimum height=1.2cm, minimum width=2.35cm, font=\scriptsize, fill=black!4}
]
\node[c] (c1) {Grounding\\cited IDs $\in$ set};
\node[c, right=0.3cm of c1] (c2) {Topology\\sensitivity};
\node[c, right=0.3cm of c2] (c3) {Corpus\\fidelity};
\node[c, right=0.3cm of c3] (c4) {Operational\\usefulness};
\node[c, right=0.3cm of c4] (c5) {Agent\\safety};
\end{tikzpicture}
\caption{Evaluation criteria checklist for twin-grounded language systems. Prefer these lenses over one paper's bar chart.}
\label{fig:ch8-eval-criteria}
\end{figure}

\begin{itemize}[leftmargin=*]
\item \textbf{Grounding:} are cited IDs in the evidence set?~\cite{ji2023,shuster2021}
\item \textbf{Topology sensitivity:} do propagation questions improve when neighbours are included?
\item \textbf{Corpus fidelity:} for codes/manuals, does RAG beat parametric recall?~\cite{lee2024ragsafety,joffe2025buildingcodes}
\item \textbf{Operational usefulness:} can outputs drive schedules (gondola staging, revisit planning)?
\item \textbf{Agent safety (emerging):} tool success under approval and geofence constraints~\cite{luo2025llmdtrobot,rawat2025llmmas}.
\end{itemize}

\subsection{Case-study snapshot (illustrative)}
Published community evaluations already show that structured retrieval helps: hybrid defect-management RAG~\cite{jeon2025hybridrag}, safety/code RAG~\cite{lee2024ragsafety,joffe2025buildingcodes}, and facility-document RAG~\cite{krutli2025fmrags}. Separately, ongoing group experiments on logistics-warehouse facades (not yet peer-published; excluded from the reference list) report illustrative multi-hop gains when asset passports and neighbour expansion are used---useful for teaching \emph{why structure matters}, not as universal constants. The following figures are therefore labeled as \emph{illustrative case sketches}, not as the chapter's scientific spine.

\begin{figure}[t]
\centering
\includegraphics[width=0.78\textwidth]{deterministic_ablation_exact_hit.pdf}
\caption{Illustrative case-study ablation sketch (optional reading): removing asset structure or neighbour expansion hurts multi-hop questions most.}
\label{fig:ch8-ablation}
\end{figure}

\begin{figure}[t]
\centering
\includegraphics[width=0.78\textwidth]{deterministic_external_exact_hit.pdf}
\caption{Illustrative case-study comparison against common RAG baselines under a shared observation-only protocol.}
\label{fig:ch8-external}
\end{figure}

\begin{figure}[t]
\centering
\includegraphics[width=0.75\textwidth]{hallucination_rate.pdf}
\caption{Illustrative case-study view of unsupported claims versus grounded configurations.}
\label{fig:ch8-halluc}
\end{figure}

\begin{figure}[t]
\centering
\includegraphics[width=0.75\textwidth]{deterministic_ksweep_coverage.pdf}
\caption{Illustrative $k$-hop sweep: coverage of propagation questions versus neighbour radius. Use to teach the operational meaning of $k$, not as a universal optimum.}
\label{fig:ch8-ksweep}
\end{figure}

\subsection{Towards predictive and cross-domain use}
Priority scoring and operational planning convert diagnosis into schedules. Transfer ideas (renovation sequencing, construction monitoring) appear in DT/agent programmes~\cite{rawat2025llmmas,luo2025llmdtrobot,lee2025multiagentvbm} and should be taught as \emph{schema reuse}, not as guaranteed forecasting. True deterioration prediction still needs temporal models beyond a single inspection epoch.

\begin{figure}[t]
\centering
\includegraphics[width=0.82\textwidth]{framework_generalizability_concept.pdf}
\caption{Conceptual transfer of twin-language patterns across O\&M tasks.}
\label{fig:ch8-transfer}
\end{figure}

\subsection{Open issues}
Topology completeness and Asset-ID hygiene dominate error modes. Standard agent benchmarks for inspection (safe tool traces) are still missing. Stronger VLM grounding on twin entities, fully VLM-native twins, richer long-term agent memory, and swarm-continuous updates belong to Chapter~9's frontier outlook, building on planning chapters for multi-UAV operations. Privacy and governance questions for building LLMs~\cite{saka2024gptconstruction,ma2026tenquestions} remain open as portfolios scale.

\paragraph{Chapter takeaways.}
\begin{enumerate}[leftmargin=*]
\item Chapter~8 is the book's decision layer: it consumes Ch.~4--7 products and answers with grounded language, not with geometry alone.
\item LLM/VLM foundations (Transformers, instruction following, multimodal interfaces) explain capability; hallucination and ID errors explain why grounding is mandatory~\cite{vaswani2017attention,brown2020,ji2023}.
\item Memory must be layered: weights $\neq$ context window $\neq$ RAG index $\neq$ twin store~\cite{lewis2020,krutli2025fmrags,pan2024b}.
\item Observation records, asset passports, and $k$-hop topology give retrieval something reliable to fetch; cite-or-abstain keeps generation honest.
\item RAG, tools/agents, and structured asset knowledge are complementary; AEC already uses them for BIM, codes, FM, and safety---facade twins add topology and Asset-ID discipline~\cite{lewis2020,zheng2023bimassistant,lee2024ragsafety,luo2025llmdtrobot,rawat2025llmmas}.
\item Reports, field identity hygiene, and approval gates are as important as model choice for decision support.
\end{enumerate}
